\documentclass[11pt,letterpaper,landscape]{article}
\usepackage[margin=.7in]{geometry}
\usepackage[T1]{fontenc}
\usepackage{lmodern,microtype,amsmath,amssymb,booktabs,graphicx,natbib,xurl,hyperref}
\hypersetup{colorlinks=true,linkcolor=blue,urlcolor=blue,citecolor=blue}
\setlength{\parindent}{0pt}\setlength{\parskip}{.5em}
\providecommand{\reportbuilddatetime}{unknown}
\input{report-build-info.tex}
\title{Geodesic MDS followed by edge-KK\\Radius, neighborhood, sampling, and scale sensitivity}
\author{\small Source: \path{~/current_projects/grip/tools/experiments/mds-edge-kk-radius/report.tex}}
\date{Build \reportbuilddatetime}
\makeatletter
\renewcommand{\maketitle}{\begin{center}{\Large\@title\par}\vspace{.5em}{\@author\par}\vspace{.3em}{\@date\par}\end{center}}
\makeatother
\newcommand{\figpage}[3]{\clearpage\section*{#2}\begin{center}\includegraphics[width=.98\textwidth,height=.70\textheight,keepaspectratio]{../figures/#1.pdf}\end{center}\small #3\normalsize}
\begin{document}
\maketitle
\section*{Question and findings}
Does refining an MDS embedding with local edge constraints restore the original
paraboloid or saddle? How does that answer depend on the neighborhood size,
the distance used to build the graph, and the optimizer's scale policy?
We examine $(x,y,x^2+y^2)$ and $(x,y,x^2-y^2)$ on disks of increasing radius.
Smooth-surface geodesics define one graph regime; ambient Euclidean chords
define the other. All fitted configurations have three coordinates.

\input{report-findings.tex}

These are paired finite-sample experiments. Increasing radius at fixed sample
size does not maintain a resolved surface mesh. Neither a thin plot nor a small
edge objective establishes a limiting dimension or recovery theorem.

\clearpage
\section*{Experimental design}
We use radii $r=1,2,4,8,16,32,64$, neighborhood sizes
$k=4,8,16,32,64,128,239$, and three independent samples of $n=240$ points.
The random seeds are 20260905--20260907. Within a sample, the radial quantiles
and angles are coupled across radii and surfaces. The first sample reuses the
audited original numerical geodesics. A nested $n=480$ sample at $r=64$ adds
240 independent points, retaining the original indices, and uses
$k=8,16,64,128,479$. Thus $8$ and $64$ test fixed $k$; $16$ and $128$ match
the original $k/n$ at $8$ and $64$. Complete graphs are explanatory controls.
This single nested size comparison does not estimate a large-$n$ limit.

Both sampling measures are included. For uniform $U\in(0,1)$ and independent
uniform angle, base-disk sampling uses $\rho=r\sqrt U$. Surface-area sampling
uses
\[
 \rho=\frac12\left(\{1+U[(1+4r^2)^{3/2}-1]\}^{2/3}-1\right)^{1/2}.
\]
Both graph surfaces have area element $\sqrt{1+4\rho^2}\,\rho\,d\rho\,d\theta$,
which gives this inverse radial distribution by integration.

An undirected edge is included if either endpoint selects the other among its
$k$ nearest neighbors; ties are broken by vertex index. In the geodesic regime,
both ranking and edge targets use the smooth-geodesic matrix $\Delta$. In the
ambient regime, both use original-coordinate chords. Disconnected unions are
explicitly augmented by a minimum spanning tree under the same distances.
No vertices are removed. Connectivity, augmentation, degree, edge lengths,
and parameter-plane radial/angular separations are recorded.

Classical MDS and stress MDS use strict graph shortest distances $D_G$.
A third initializer in the geodesic regime is classical scaling of full
$\Delta$, held fixed across $k$. Stress MDS uses the grip adapter to
\texttt{smacof::mds(type="ratio")}: stress majorization with ratio disparities
and internal normalization \citep{mair2022}. The adapter restores the
raw-stress-optimal coordinate scale. Each graph has one classical and two
Gaussian starts, at most 1,000 iterations per start and tolerance $10^{-8}$.
Gaussian draws are matched across $k$; the adapter optimally rescales each
start against its own target matrix. The least achieved raw stress selects
the fit; global optimality is not claimed.

\clearpage
\section*{The pilot changed the scale controls}
Every target and initializer is first divided by the RMS full smooth-geodesic
distance. The declared edge-KK workflow uses density-to-uniform stiffness
$(0,.25,.5,.75,1)$ with 200 steps per stage and a profiled target scale.
A uniform-only run has the same total allowance of 1,000 steps. The original
72-graph pilot exposed severe contraction in some geodesic fits. We retained
those outputs and repeated the pilot with additional fixed-target-scale runs,
using both schedules. Those controls were then applied throughout the study.
The package defaults were not changed.

Write $\ell_e>0$ for edge targets, $d_e(Z)=\|z_i-z_j\|$ for embedded edge
lengths, and $\kappa_e\geq0$ for the current stiffness (nonnegative, with a positive weighted target norm). The implemented objective is
\[
 E(Z,a)=\frac12\sum_e\kappa_e[d_e(Z)-a\ell_e]^2,
 \qquad a(Z)=\frac{\sum_e\kappa_e d_e(Z)\ell_e}{\sum_e\kappa_e\ell_e^2}.
\]
Since $d_e(cZ)=c\,d_e(Z)$, for $c>0$ we have $a(cZ)=c\,a(Z)$ and
\[
 \boxed{E(cZ,a(cZ))=c^2E(Z,a(Z)).}
\]
Thus uniform contraction reduces any positive profiled loss without changing
shape or relative edge errors. A vanishing absolute objective or gradient is
not a certificate of shape convergence. This statement follows directly from
the implemented loss and is independently checked on saved pilot fits.
It does not invalidate relative scores computed at a common calibrated scale.
A scale-invariant alternative would divide the minimized numerator by
$\sum_e\kappa_e d_e^2$, giving
\[
 Q(Z)=1-\frac{(\sum_e\kappa_e d_e\ell_e)^2}
 {(\sum_e\kappa_e d_e^2)(\sum_e\kappa_e\ell_e^2)}.
\]
When $\sum_e\kappa_e d_e^2>0$, this also equals raw edge stress minimized
over a coordinate multiplier, divided by the weighted target norm squared.
The identity follows by minimizing a scalar quadratic. It explains why fixing
target scale is a valid control; it does not make unnormalized profiled loss
scale invariant or equate the finite-budget trajectories.

Fixed-scale controls use $a=1$, so contraction cannot remove a positive target.
For a complete graph with uniform stiffness, $2E(Z,1)$ equals all-pair raw
distance stress when the same targets are used. Tiny numerical violations of
triangle inequalities in smooth inputs can distinguish direct targets from
strict graph distances; these discrepancies are measured separately.
For complete ambient graphs, the original coordinates exactly realize the targets.

For shape displays and end-to-end path scoring, let
$b=\sum_e d_e\ell_e/\sum_e\ell_e^2$. We divide fitted coordinates or path
lengths by this edge-derived $b$, then restore physical units. No smooth-reference
scale is fitted. Saved coordinates also retain the optimizer's returned scale.
Plots use equal spatial units and at most one additional uniform $r^2$ divisor.

\clearpage
\section*{What the diagnostics mean}
For observed length vector $y$ and positive target vector $t$, define
$a=y^\top t/(t^\top t)$. Edge and fixed-path relative errors are
$\|y-at\|/\|at\|$, each with its own diagnostic scale. Every graph uses all
unordered vertex pairs and the same retained routes for every candidate.
The graph constructor's near-tie route convention can produce slightly longer
routes than strict $D_G$; MDS receives symmetric strict distances and route
scoring uses the retained route lengths. These roles are not interchanged.

Chord raw stress is $\sum_{i<j}(d_{ij}-D_{G,ij})^2$ at returned physical scale.
Profiled chord Stress-1 is $\|d-aD_G\|/\|d\|$. These are different statistics;
raw stress after a contracting profiled refinement is not its optimized criterion.
Graph/reference error is $\|D_G-\Delta\|/\|\Delta\|$ in physical units.
End-to-end path/reference error is $\|L_Z/b-\Delta\|/\|\Delta\|$, where
$L_Z$ sums embedded edge lengths along retained routes. Small graph-to-layout
error cannot substitute for graph-to-surface validation.

Coordinate error is relative Frobenius error after similarity alignment to
known generating coordinates with known point correspondence. It is not a
closest-surface distance or a test of topology. Let
$s_1\geq s_2\geq s_3$ be singular values of the centered coordinates.
A line limit requires $s_2/s_1\to0$; a nondegenerate plane requires
$s_3/s_2\to0$ while $s_2/s_1$ stays nonzero. Both ratios are needed.
We report finite-radius ratios and all three dimensional singular values.

A useful deterministic bound explains one implication of edge agreement.
If every retained edge satisfies $|d_e/(b\ell_e)-1|\leq\epsilon$, then every
retained path $\gamma$ satisfies
\[
 \left|\frac{L_Z(\gamma)}b-L_G(\gamma)\right|
 \leq\sum_{e\in\gamma}|d_e/b-\ell_e|
 \leq\epsilon\sum_{e\in\gamma}\ell_e
 =\epsilon L_G(\gamma).
\]
This is a graph-to-layout statement. Small mean edge stress does not establish
its maximum-relative-error hypothesis, and the bound does not identify a unique
surface realization. A one-dimensional intrinsic metric can also require a
higher-dimensional Euclidean approximation: intrinsic dimension and Euclidean
span answer different questions. For example, four leaves at distance $t>0$
from the center of a metric star have all interleaf distances $2t$. Their
classical Gram matrix is $2t^2J$, where $J=I-\mathbf1\mathbf1^\top/4$ has rank
three. Thus even this four-point subset of an intrinsically one-dimensional
tree requires three Euclidean coordinates for exact distance realization. The existing direct-geodesic MDS line/tree
analyses therefore do not automatically determine edge-KK's limit.

\clearpage
\section*{Numerical reliability and optimizer sensitivity}
\input{report-validation.tex}

At $r=64$, replicate 1, $k=8,32,128,239$, additional fits start from tiny
three-dimensional perturbations, independent random coordinates, and original
coordinates. Perturbations have entrywise RMS $10^{-4}$ in normalized units.
Separate continuations add 2,000 uniform-stiffness steps to each selected
classical- and stress-initialized edge-KK result. Both scale policies are
included. Random and original-coordinate configurations are shared across the
classical and stress branch comparisons; their duplicate table rows are not
independent replications. These exploratory controls do not replace primary fits.
A planar start can have zero transverse gradient for a differentiable
distance-based objective; a perturbation check addresses particular fits,
not every planar global minimizer.

\input{report-optimizer.tex}

\clearpage
\section*{Neighborhood, sampling, and sample size}
\input{report-sensitivity.tex}

\clearpage
\section*{Implications for the grip software article}
\input{report-implications.tex}

The full radius, $k$, sampling, and size study belongs in a supplement or a
separate research report. A concise main-text comparison should distinguish
classical scaling, stress MDS, and edge-only refinement; explain the graph
regime; and avoid equating path fidelity with generating-surface recovery.
The scale issue needs explicit treatment before using loss or stopping values
as evidence for the profiled workflow. This report records an experiment and
its controls; integration into the main manuscript is a later phase.

\input{report-figures.tex}
\clearpage
\section*{Reproducibility and references}
The accompanying protocol, R/Python sources, compact distance matrices,
coordinates, full stratified tables, input/source hashes, and validation
records form the reproducibility bundle. Data-generation and fitting caches
are separate from the compact artifacts used to render this report.
No build or validation step requires private agent notes.
\bibliographystyle{plainnat}
\bibliography{../references}
\end{document}
